\documentclass[11pt]{article}
% ---------------------------------------------------------------------------
%  THE MODEL INPUT, FEATURE BY FEATURE
%  Detailed definitions (formulas), preparation recipe, and the physical and
%  empirical justification of every input the ML model receives.
%  Companion to gme_ml_physics_overview.tex.
%  Compile:  pdflatex gme_ml_input_features.tex   (twice for TOC)
% ---------------------------------------------------------------------------
\usepackage[a4paper,margin=2.2cm]{geometry}
\usepackage{amsmath,amssymb}
\usepackage{booktabs,array}
\usepackage{xcolor}
\usepackage{enumitem}
\usepackage[hidelinks]{hyperref}

\newcommand{\CGME}{\ensuremath{C_{\mathrm{GME}}}}
\newcommand{\rdm}{\ensuremath{\Gamma^{(n)}}}
\newcommand{\Bb}{\ensuremath{B_b}}
\newcommand{\Bt}{\ensuremath{\widetilde B_b}}

\title{\textbf{The Model Input, Feature by Feature}\\[4pt]
       \large What the GME-estimation model receives, how every number is
       computed from the reduced density matrix, and why it is necessary}
\author{Hubbard-GME project}
\date{\today}

\begin{document}
\maketitle
\tableofcontents
\bigskip

% ===========================================================================
\section{The problem the input must solve}
% ===========================================================================
The model must predict the genuine multipartite entanglement
\[
  \CGME(\rho)\;=\;\min_{\{p_i,\phi_i\}}\;\sum_i p_i\,C(\phi_i),
  \qquad \rho=\sum_i p_i\,|\phi_i\rangle\langle\phi_i|,
\]
--- an expensive convex-roof optimisation (seconds to hours per state) ---
from \emph{cheap} numbers computable from $\rho$ in microseconds. The input
design must satisfy four hard constraints:
\begin{enumerate}[nosep]
  \item \textbf{Computable from $\rho$ alone.} No Hamiltonian parameters, no
        temperature, no solver output --- otherwise the model cannot transfer
        to molecules or leak-free evaluation is impossible.
  \item \textbf{Fixed length for every cluster size} $n=2\ldots10$, so one
        model serves all orders (a training table must be rectangular).
  \item \textbf{Scalable}: the length must grow at most linearly with the
        design cap, while $\rho$ itself grows as $4^n$ (at $n{=}10$:
        $91{,}866$ unique off-diagonal entries --- unusable as input).
  \item \textbf{Convention-robust}: invariant under the orbital sign/phase
        ambiguities of the construction (verified; Section~\ref{sec:gauge}).
\end{enumerate}

% ===========================================================================
\section{How the dataset is prepared}
% ===========================================================================
\subsection{The reduced density matrix (the raw material)}
For an eigenstate $|\Psi_m\rangle$ of the system (Hubbard ring of $N$ sites,
or a molecular active space) and a cluster of $n$ orbitals, the $n$-spin
reduced density matrix is built element by element,
\[
  \rdm_{\sigma'\sigma}
  \;=\;\Big\langle\Psi_m\Big|\,\prod_{k=1}^{n}
     d^{\dagger}_{k,\sigma'_k}\,d^{\phantom\dagger}_{k,\sigma_k}
     \Big|\Psi_m\Big\rangle ,
  \qquad \sigma,\sigma'\in\{\uparrow,\downarrow\}^{n},
\]
then normalised to unit trace. It is a $2^n\times 2^n$ matrix over the
cluster's spin configurations: the diagonal holds the \emph{populations}
(probabilities of each spin pattern), the off-diagonal holds the
\emph{coherences} (quantum superposition amplitudes between patterns). The
construction projects onto single occupation of every cluster site, which
separates spin from charge and is the origin of the gauge robustness proven
in Section~\ref{sec:gauge}.

\subsection{Thermal states}
Because the per-eigenstate RDMs do not depend on temperature, a thermal
(Gibbs) state at temperature $T$ costs no new diagonalisation:
\[
  \rdm(T)\;=\;\frac{\sum_m e^{-(E_m-E_0)/T}\;\rdm[\Psi_m]}
                   {\sum_m e^{-(E_m-E_0)/T}} .
\]
The dataset uses the grid $T\in\{0,\,0.1,\,0.25,\,0.5,\,1,\,2\}$; the hot
points supply the essential $\CGME\approx0$ examples (class balance).

\subsection{The label}
Every record is labelled by the GCG convex-roof solver:
$C_{\mathrm{target}}=C_{\mathrm{gcg}}$, computed with a fixed budget
(8 restarts, 150 iterations, $k_{\max}$ rank $+2$). The label's own
numerical uncertainty is being measured by a repeated-ensemble experiment;
features can only be judged against errors larger than that floor.

\subsection{One record}
One record $=$ one $(\text{system},\text{state or }T,\text{cluster})$ triple:
the full $\rho$ (archived), the configuration metadata (never model input),
the label, and the feature vector defined below. Records are grouped by a
hash of the physical configuration; train/test splits are made \emph{by
group}, so the model is never tested on a Hamiltonian it trained on.

% ===========================================================================
\section{The $S_z$ block structure (the organising principle)}
% ===========================================================================
Every parent state conserves total $S_z$, so the RDM is \emph{exactly}
block-diagonal in the cluster's total $S_z$ (verified to $10^{-15}$ over the
full dataset, both symmetry sectors). Index the blocks by the number of down
spins $b$:
\[
  \rho \;=\; \bigoplus_{b=0}^{n} \Bb ,
  \qquad
  \dim \Bb = \binom{n}{b},
  \qquad
  S_z^{(b)} = \frac{n-2b}{2} .
\]
Example $n=3$: blocks of size $1,3,3,1$ ($S_z=+\tfrac32,+\tfrac12,
-\tfrac12,-\tfrac32$). Everything outside the blocks is structurally zero
and carries no information; the features therefore summarise each block
separately and never waste dimensions on entries known to vanish.

Two conventions apply to \emph{all} features below:
\begin{itemize}[nosep]
  \item \textbf{Strict upper triangle}: off-diagonal sums run over $i<j$
        within a block, so no coherence is counted twice (Hermiticity makes
        $(i,j)$ and $(j,i)$ redundant).
  \item \textbf{Structural zeros excluded}: statistics (max, fractions,
        quantiles) are taken only over valid within-block entries.
\end{itemize}
The RDMs are real to machine precision (verified: $\max|\mathrm{Im}\,\rho|=0$
across all records), so real-part formulas are exact, not approximations.

% ===========================================================================
\section{Feature-by-feature definitions}\label{sec:features}
% ===========================================================================
The vector reserves $11$ block slots ($b=0\ldots10$, design cap $n\le10$);
for a record of cluster size $n$ only slots $b\le n$ are ``real''. Features
come in three groups: \emph{per-block} (repeated for each slot),
\emph{global}, and the \emph{size variable}.

\subsection{Cluster size $n$}
\[
  \fbox{$n$}\qquad\text{(1 number)}
\]
Sets the dimension $2^n$, the block layout $\binom{n}{b}$, and the
entanglement order being asked about ($n$-partite). Without it the model
cannot interpret the block pattern. \emph{Why necessary:} the same block
values mean different physics at different $n$ (e.g.\ a coherence of $0.5$
is maximal for $n=2$ but moderate for $n=6$).

\subsection{Block mask $m_b$}
\[
  m_b=\begin{cases}1,& b\le n\\ 0,& b>n\end{cases}
\]
Distinguishes a \emph{padded} slot (does not exist for this $n$) from a
\emph{real but empty} block ($p_b=0$ because the state has no weight in that
$S_z$ sector --- physically meaningful, e.g.\ $S_z{=}0$-sector states never
populate the extreme blocks). Without the mask these two very different
situations look identical to the model.

\subsection{Block population $p_b$}
\[
  p_b=\operatorname{Tr}\Bb\;\in[0,1],\qquad \sum_b p_b=1 .
\]
\emph{Meaning:} the probability that the cluster is found with $b$ down
spins ($S_z=(n-2b)/2$). \emph{Why necessary:} the population profile encodes
the spin composition of the state --- ground states at half filling
concentrate weight near $S_z\!\approx\!0$; heating in the full sector pushes
weight into the outer (ferromagnetic-like) blocks; a $\CGME$ model must know
\emph{where} the state lives before asking how coherent it is there.

\subsection{The conditional block $\Bt$ (not a feature --- a normalisation)}
\[
  \Bt=\frac{\Bb}{p_b}\qquad (p_b>0),\qquad \operatorname{Tr}\Bt=1 .
\]
All internal-structure features below are computed from $\Bt$, not $\Bb$.
\emph{Why:} it disentangles two independent facts --- \emph{how much}
probability sits in a sector ($p_b$) and \emph{how organised} that
probability is inside the sector ($\Bt$). Without the normalisation, a
weakly-populated but perfectly coherent block would look identical to a
well-populated but incoherent one.

\subsection{Conditional block purity $q_b$}
\[
  q_b=\operatorname{Tr}\!\big(\Bt^{\,2}\big)\;\in\Big[\tfrac{1}{\dim\Bb},\,1\Big].
\]
\emph{Meaning:} how mixed the sector is internally: $q_b=1$ --- the sector
holds a single pure superposition; $q_b=1/\dim\Bb$ --- maximally mixed
(classical noise). \emph{Why necessary:} mixedness is precisely what makes
the convex roof nontrivial --- it measures how much room the ``least
entangled explanation'' has to explain the correlations classically.
Thermal records are separated from ground-state records mainly by this
feature (which is why the temperature $T$ itself is \emph{not} needed as an
input).

\subsection{Coherence strength $C_{2,b}$}
\[
  C_{2,b}=\Big(\,2\!\!\sum_{i<j\in b}\!\big|(\Bt)_{ij}\big|^2\Big)^{1/2}.
\]
\emph{Meaning:} the total magnitude of quantum superposition inside the
sector --- the single most entanglement-relevant number. \emph{Why
necessary:} entanglement lives in the off-diagonals. The canonical
demonstration is the three-spin W state vs.\ its classical mixture:
\[
  |W\rangle=\tfrac{1}{\sqrt3}\big(|{\downarrow\uparrow\uparrow}\rangle
  +|{\uparrow\downarrow\uparrow}\rangle
  +|{\uparrow\uparrow\downarrow}\rangle\big)
  \quad\text{vs.}\quad
  \rho_{\mathrm{mix}}=\tfrac13\!\sum_{\text{same 3}}|{\cdot}\rangle\langle{\cdot}| .
\]
Both have identical diagonals $(\tfrac13,\tfrac13,\tfrac13)$ and identical
populations; they differ \emph{only} in $C_{2,b}$ ($0.816$ vs $0$) and $q_b$
($1$ vs $\tfrac13$). A diagonal-only model cannot tell an entangled state
from classical noise --- and empirically collapses under extrapolation
($R^2=0.08$ on hold-out-$L$).

\subsection{Signed coherences $s^{+}_b$, $s^{-}_b$}
\[
  s^{+}_b=\sum_{i<j\in b}\max\!\big((\Bt)_{ij},\,0\big),
  \qquad
  s^{-}_b=\sum_{i<j\in b}\max\!\big(-(\Bt)_{ij},\,0\big).
\]
\emph{Meaning:} the total weight of positive vs negative coherences. In our
fixed collinear-$S_z$ computational basis the \textbf{sign carries spin
physics}: negative coherences are singlet-type (antisymmetric) couplings ---
the antiferromagnetic fingerprint of the Hubbard ground state --- while
positive coherences are triplet-type (symmetric) couplings. \emph{Why two
numbers instead of one signed sum:} a mixed pattern $(+a,-a)$ and no
coherence at all both give a signed sum of zero; the pair
$(s^{+},s^{-})$ keeps them distinct. \emph{Redundancy rule:} the signed sum
and the $L_1$ coherence are linear combinations,
\[
  s_b=s^{+}_b-s^{-}_b,\qquad L_{1,b}=s^{+}_b+s^{-}_b ,
\]
so $(s^{+}_b,s^{-}_b)$ is stored and $(s_b,L_{1,b})$ is \emph{not} --- for
our (real) RDMs they contain exactly the same information.

\emph{Convention-safety:} these signed features are invariant under
per-orbital sign/phase changes --- proven algebraically (each factor
$d^{\dagger}_{k\sigma'}d_{k\sigma}$ carries $e^{+i\theta_k}e^{-i\theta_k}=1$)
and verified numerically on both the Hubbard ring
($\max|\Delta\rho|=4.6\times10^{-15}$) and benzene PM orbitals
($1.2\times10^{-13}$); see Section~\ref{sec:gauge}.

\subsection{Hamming-resolved coherence distribution $\bar R_{b,d}$}
Group the coherences of a block by the Hamming distance between the two spin
configurations they connect (the number of sites whose spin differs):
\[
  H_{b,d}=\sum_{\substack{i<j\in b\\ \mathrm{Ham}(i,j)=d}}
          \big|(\Bt)_{ij}\big|^{2},
  \qquad d=2,4,6,8,10 ,
\]
(within a fixed-$S_z$ block only \emph{even} distances occur --- flipping one
spin up requires flipping another down). The stored feature is the
\emph{normalised distribution},
\[
  \bar R_{b,d}=\frac{H_{b,d}}{\sum_{d'}H_{b,d'}} ,
\]
because the total is already carried by $C_{2,b}$
($\sum_d H_{b,d}=C_{2,b}^2/2$) --- storing $C_2$, $H$ and $R$ together would
be redundant.

\emph{Meaning:} \emph{where} the coherence connects: $d=2$ --- local spin
exchanges (one pair swapped, e.g.\ nearest-neighbour singlet motion);
$d=n$ (max) --- superpositions of \emph{globally} different spin patterns
(N\'eel$\,\leftrightarrow\,$anti-N\'eel), the genuinely multipartite,
GHZ-like coherences. \emph{Why necessary:} the summary $C_{2,b}$ says how
much coherence exists, $\bar R_{b,d}$ says whether it is pairwise-local or
many-body-global --- precisely the distinction between a state entangled
only in pairs and one with true multipartite structure. This is the compact,
scalable substitute for the raw coherence \emph{locations} that the
$+0.07$ hold-out-$R^2$ advantage of raw entries suggests matters.
\emph{Caution:} Hamming distance is configuration distance, \emph{not}
spatial distance between orbitals.

\subsection{Strongest coherence $C_{\max,b}$ (shape family)}
\[
  C_{\max,b}=\max_{i<j\in b}\big|(\Bt)_{ij}\big| .
\]
\emph{Meaning/why:} distinguishes \emph{concentrated} coherence (one
dominant pair --- Bell-like) from \emph{distributed} coherence (many small
terms --- W-like) at equal $C_{2,b}$.

\subsection{Significant-coherence fraction $f_{\tau,b}$ (shape family)}
\[
  f_{\tau,b}=\frac{\#\{(i,j):\,i<j,\;|(\Bt)_{ij}|>\tau\}}
                  {\binom{\dim\Bb}{2}},
  \qquad
  \boxed{\tau=10^{-6}\ \text{(documented, absolute, on }\Bt\text{)}}
\]
\emph{Meaning/why:} the sparsity of the coherence pattern --- a second,
complementary concentration measure. The threshold is fixed and documented
so the feature is reproducible; entries below $\tau$ are numerically
indistinguishable from zero at solver precision.

\subsection{Global summaries (4 numbers)}
\[
  C_2^{\mathrm{glob}}=\Big(2\!\!\sum_{i<j}|\rho_{ij}|^2\Big)^{1/2},\quad
  \max_{i<j}|\rho_{ij}|,\quad
  \min_{i<j}\rho_{ij},\quad
  \max_{i<j}\rho_{ij},
\]
computed over all valid within-block positions of the \emph{unnormalised}
$\rho$. \emph{Why:} a whole-matrix fingerprint that survives even when
individual blocks are nearly empty; the signed extrema flag the single
strongest singlet and triplet coupling in the state.

% ===========================================================================
\section{The three nested vectors (exact counts)}
% ===========================================================================
\begin{center}
\begin{tabular}{@{}llcc@{}}
\toprule
Version & Per-block content ($\times11$ slots) & Formula & Total \\
\midrule
V1 core     & $m_b,\;p_b,\;q_b,\;C_{2,b},\;s^{+}_b,\;s^{-}_b$
            & $11\times6+1+4$ & \textbf{71} \\
V2 +Hamming & V1 $+\;\bar R_{b,2},\bar R_{b,4},\bar R_{b,6},
              \bar R_{b,8},\bar R_{b,10}$
            & $71+11\times5$ & \textbf{126} \\
V3 +shape   & V2 $+\;C_{\max,b},\;f_{\tau,b}$
            & $126+11\times2$ & \textbf{148} \\
\bottomrule
\end{tabular}
\end{center}
Padded slots ($b>n$): all entries $0$ and $m_b=0$. Invalid Hamming distances
($d>n$): $0$. The versions are nested so the ablation isolates each family's
contribution; the final production vector is frozen by experiment (hard-band
accuracy, hold-out-$L$, molecular transfer), not by argument.

% ===========================================================================
\section{Why not the raw matrix? (the evidence)}\label{sec:evidence}
% ===========================================================================
Six input designs, gradient-boosted trees, split grouped by Hamiltonian
(4{,}339 records; hard band $=|C-0.01|<0.05$):
\begin{center}
\small
\begin{tabular}{@{}lcccc@{}}
\toprule
Input & MAE & $R^2$ & hard-band acc. & hold-out-$L{=}6$ $R^2$ \\
\midrule
diagonals only          & 0.0563 & 0.902 & 0.812 & 0.078 \\
raw off-diagonals only  & 0.0232 & 0.982 & 0.851 & 0.517 \\
33 block summaries      & 0.0087 & 0.997 & 0.928 & 0.790 \\
partial enriched        & 0.0076 & 0.998 & 0.918 & 0.799 \\
raw $+$ summaries       & \textbf{0.0074} & \textbf{0.998} & 0.928
                        & \textbf{0.864} \\
\bottomrule
\end{tabular}
\end{center}
\begin{enumerate}[nosep]
  \item \textbf{Populations alone fail} (hold-out $R^2=0.08$): coherence is
        not optional --- it is where entanglement lives.
  \item \textbf{Summaries beat raw entries used alone} ($0.0087$ vs
        $0.0232$): physics-shaped compression is worth more than raw data at
        this dataset size.
  \item \textbf{Raw entries add extrapolation value on top} ($0.864$ vs
        $0.790$): location information matters; the scalable way to capture
        it is $\bar R_{b,d}$ (V2), not $91{,}866$ raw numbers.
\end{enumerate}

% ===========================================================================
\section{Convention-robustness (gauge) and its limits}\label{sec:gauge}
% ===========================================================================
Under a per-orbital phase change $\varphi_k\to e^{i\theta_k}\varphi_k$
(sign flips $\theta_k=\pi$ included --- the ambiguity of localized
molecular orbitals), each RDM factor transforms as
\[
  d^{\dagger}_{k\sigma'}d^{\phantom\dagger}_{k\sigma}
  \;\longrightarrow\;
  e^{+i\theta_k}e^{-i\theta_k}\,
  d^{\dagger}_{k\sigma'}d^{\phantom\dagger}_{k\sigma}
  \;=\;
  d^{\dagger}_{k\sigma'}d^{\phantom\dagger}_{k\sigma},
\]
so \rdm{} --- and every feature above, including the signed ones --- is
\textbf{exactly invariant}. Verified numerically: Hubbard ring
$\max|\Delta\rho|=4.6\times10^{-15}$; benzene CAS(6,6) with PM orbitals
$1.2\times10^{-13}$.

\emph{Limits (stated honestly):} invariance holds for a common phase on the
spatial orbital shared by both spins. It does \emph{not} cover local
spin-basis rotations (which redefine the computational basis) or general
orbital \emph{rotations} --- the latter change the physical partition and
\emph{should} change the answer (canonical vs localized benzene is real
physics, not an artifact). Separately, localized orbitals must be
\emph{canonicalised} (deterministic ordering + sign convention) before
dataset generation, because localization can return the same orbitals in a
different order --- a bookkeeping ambiguity that relabels cluster indices.

% ===========================================================================
\section{What is deliberately NOT an input}
% ===========================================================================
\begin{center}
\begin{tabular}{@{}lp{8.6cm}@{}}
\toprule
Excluded & Reason \\
\midrule
$L,\,t,\,U,\,N_e$, topology & ablation: adds nothing once $\rho$ is
  described; excluding them makes the model Hamiltonian-agnostic
  (molecular transfer) \\
temperature $T$ & readable from the block purities; excluding it lets one
  model serve ground, excited and thermal states \\
full $\rho$ entries & $4^n$ scaling breaks the single-model design;
  archived in every record for audits and Phase-2 \\
GCG output (bounds, histories, decompositions) & \textbf{target leakage}:
  at prediction time no GCG output exists --- producing it is the very cost
  the model removes \\
eigenvalue-only (spectral) features & unitary-invariant, hence blind to
  coherences (a GHZ state and a product state can share a spectrum); kept
  only as a legacy comparison set \\
\bottomrule
\end{tabular}
\end{center}

% ===========================================================================
\section{Summary}
% ===========================================================================
Each record: $\rho$ (archived) $\to$ $S_z$ blocks $\to$ per-block
$(m_b,p_b,q_b,C_{2,b},s^{+}_b,s^{-}_b)$ $+$ optional
$\bar R_{b,d}$ and shape features $+$ $n$ $+$ 4 global summaries $\to$ a
fixed-length, gauge-robust vector (71/126/148) $\to$ model $\to$
$\widehat C_{\mathrm{GME}}$, trained against GCG labels under
group-split evaluation. Every feature has a formula, a physical meaning, and
an empirical reason to exist; everything the features cannot see is still
archived in the record, so the design can be revisited without recomputing
a single label.
\end{document}
